% !TEX root = ../main.tex
\chapter{Model Systems}
\label{chapter:model_systems}

This chapter defines the physical models used in the numerical
analyses of this thesis. The open-quantum-system framework was developed in
~\autoref{chapter:oqs_open_quantum_systems}, and the spectroscopy formalism was
developed in~\autoref{chapter:spctr_spectroscopy}. Here those ingredients are
combined into one common model that is then specialised to the concrete
benchmark cases used later in~\autoref{chapter:results}.
For the monomer and dimer, the concrete parameter sets introduced below are
chosen to match benchmark cases studied
in~\cite{mancaletal2006twodimensionalopticalthreepulse,pisliakovetal2006twodimensionalopticalthreepulse}.

The chapter introduces a general
\(N_{\mathrm{sites}}\)-site model of which the monomer and dimer are
concrete \(N_{\mathrm{sites}}=1\) and \(N_{\mathrm{sites}}=2\) realisations of. 

\section{General N-Site Model}
\label{sec:model_ensemble}

In the general case, we consider a system of \(N_{\mathrm{sites}}\) coupled sites. Each site is
locally truncated to a two-level system, so every site can be either in its
ground state or singly excited.

\subsection{States, Truncation, and Rotating-Frame Bookkeeping}

For \(N_{\mathrm{sites}}>1\), the Hilbert space is further truncated at the double-excitation
manifold. The resulting basis therefore has three bands:
\begin{itemize}
\item the ground state \(\ket{g}\),
\item the single-excitation states \(\ket{n}\) with \(n\in\{1,\dots,N_{\mathrm{sites}}\}\),
\item and the double-excitation states \(\ket{mn}\) with \(1\le m<n\le N_{\mathrm{sites}}\).
\end{itemize}
Here \(\ket{n}\) means that only site \(n\) is excited, while \(\ket{mn}\)
means that sites \(m\) and \(n\) are excited. Since each site is itself only a
two-level system, a given site can never be excited twice. The
double-excitation sector therefore always consists of excitations on two
different sites.

If only the one-excitation manifold is modeled, the basis is
\begin{equation*}
\mathcal{B}_{1\mathrm{ex}} = \left\{ \ket{g}, \ket{1}, \ldots, \ket{N_{\mathrm{sites}}} \right\},
\end{equation*}
so that the Hilbert-space dimension is \(1+N_{\mathrm{sites}}\). If the two-excitation manifold
is used as well, the basis is enlarged to
\begin{equation*}
\mathcal{B}_{2\mathrm{ex}} = \left\{ \ket{g}, \ket{1}, \ldots, \ket{N_{\mathrm{sites}}}, \ket{12}, \ket{13}, \ldots \right\},
\end{equation*}
and the Hilbert-space dimension becomes
\begin{equation*}
\dim \mathcal{H} = 1 + N_{\mathrm{sites}} + \binom{N_{\mathrm{sites}}}{2}.
\end{equation*}

The carrier-rotating representation used later in the numerical propagation is
defined at model level through the excitation-number operator
\(\hat{N}\). It assigns 0 to \(\ket{g}\), 1 to the single-excitation
manifold \(\{\ket{n}\}\), and 2 to the double-excitation manifold
\(\{\ket{mn}\}\). The corresponding
rotating-frame unitary is
\begin{equation}
	U(t)=\exp\!\left(i\omega_{L} t \hat{N}\right).
	\label{eq:UnitaryRotatingFrame}
\end{equation}
The driven model is therefore obtained by coupling the field-free Hamiltonian to
the optical field through the dipole operator, with the rotating-frame and RWA
conventions inherited from~\autoref{chapter:spctr_spectroscopy}.

\subsection{Hamiltonian}

In this truncated basis, a generic Frenkel-exciton Hamiltonian reads
\begin{equation*}
\hat{H}_{0}^{(N_{\mathrm{sites}})}
=
\sum_{n=1}^{N_{\mathrm{sites}}}\omega_{n}\,\ket{n}\!\bra{n}
+
\sum_{m<n}^{N_{\mathrm{sites}}}J_{mn}\,\bigl(\ket{m}\!\bra{n}+\ket{n}\!\bra{m}\bigr)
+
\sum_{m<n}^{N_{\mathrm{sites}}}(\omega_{m}+\omega_{n})\,\ket{mn}\!\bra{mn}.
\end{equation*}
The physical structure is therefore the same across the whole model: site
energies on the diagonal, coherent couplings between sites.

\subsection{Geometry and Couplings}

The couplings \(J_{mn}\) may be chosen phenomenologically or generated from
transition-dipole interactions. For the latter,
\begin{equation*}
J_{mn} \propto \frac{1}{\lVert \boldsymbol{r}_{mn} \rVert^{3}}
\left[
\boldsymbol{\mu}_{m}\cdot\boldsymbol{\mu}_{n}
- 3\,
(\boldsymbol{\mu}_{m}\cdot\hat{\boldsymbol{r}}_{mn})
(\boldsymbol{\mu}_{n}\cdot\hat{\boldsymbol{r}}_{mn})
\right],
\end{equation*}
with site positions \(\boldsymbol{r}_{n}\) and displacements
\(\boldsymbol{r}_{mn}=\boldsymbol{r}_{m}-\boldsymbol{r}_{n}\)
~\cite{griffiths2013introductionelectrodynamics,lehmberg1970radiationatomsystem,masters2014pathsforstersresonance}.


\subsection{Dipole Operator}

The dipole operator connects only adjacent manifolds,
\begin{equation*}
\hat{\mu}^{(N_{\mathrm{sites}})}
=
\sum_{n=1}^{N_{\mathrm{sites}}}
\Bigl(
\mu_{gn}\,\ket{g}\!\bra{n}
+
\mu_{ng}\,\ket{n}\!\bra{g}
\Bigr)
+
\sum_{m<n}^{N_{\mathrm{sites}}}\sum_{\ell\in\{m,n\}}
\Bigl(
\mu_{\ell,mn}\,\ket{\ell}\!\bra{mn}
+
\mu_{mn,\ell}\,\ket{mn}\!\bra{\ell}
\Bigr),
\end{equation*}
where the second sum is only needed when the two-excitation manifold is
used. This is the structural ingredient that allows excited-state
absorption for the multilevel systems.

\subsection{System--Bath Coupling Operators}

For the microscopic open-system treatment, the system operators in the
interaction Hamiltonian~\autoref{eq:oqs_Interaction_Hamiltonian} are chosen to
act locally on the site \(i\) as
\begin{equation*}
\hat{S}_{i}
=
\ket{i}\!\bra{i}
+
\sum_{\substack{m<n\\ i\in\{m,n\}}}\ket{mn}\!\bra{mn},
\qquad i\in\{1,\dots,N_{\mathrm{sites}}\}\,,
\end{equation*}
where the second term is only present when the double-excitation manifold is
retained. Thus the bath associated with site \(i\) acts not only on the
one-excitation state \(\ket{i}\), but also on every double-excitation state
that contains site \(i\). These are the system coupling operators used later
for the microscopic Redfield model of the dimer and the larger
\(N_{\mathrm{sites}}\)-site examples.

\subsection{Static Disorder}

Whenever inhomogeneous broadening is included, the relevant electronic energies
are sampled over a static disorder ensemble before the signal is averaged. In
the general \(N_{\mathrm{sites}}\)-site model these fluctuating quantities are the site energies
\(\omega_n\). A macroscopic motivation for this construction was given
in~\autoref{sec:spctr_spectroscopy_motivation,spctr:sec:macroscopic_polarisation}.
These environments are modeled to be completely uncorrelated. How this is implemeted numerically is given in
~\autoref{app:sec:inhomogeneous_broadening}.

\section{Monomer as the \texorpdfstring{$N_{\mathrm{sites}}=1$}{Nsites=1} Limit}
\label{sec:model_tls}

The monomer is the minimal benchmark used throughout the thesis. It consists of only one transition. This gives
the simplest system in which the full 2DES workflow can be validated before
coupling, cross peaks, and multilevel pathway structure are introduced.

\subsection{Hamiltonian and Dipole Operator}

For \(N_{\mathrm{sites}}=1\), the field-free Hamiltonian and optical dipole operator reduce to
the two-level expressions already introduced in
~\cref{eq:spctr_tls_hamiltonian,eq:spctr_tls_dipole}. For later reference, they
are repeated here:
\begin{equation*}
\hat{H}_{0}^{(\mathrm{TLS})} = \omega_{eg}\,\ket{e}\!\bra{e},
\end{equation*}
and
\begin{equation*}
\hat{\mu}^{(\mathrm{TLS})}
= \mu_{ge}\,\ket{g}\!\bra{e} + \mu_{eg}\,\ket{e}\!\bra{g},
\qquad
\mu_{eg}=\mu_{ge}^{*}.
\end{equation*}

\subsection{Phenomenological Lindblad Bath}

For the monomer calibration benchmark, the open-system treatment is kept
phenomenological. In the notation of the general model,
the local excitation projector reduces to
\begin{equation*}
\hat{S}_{e} = \ket{e}\!\bra{e}.
\end{equation*}
The benchmark monomer itself is then coupled to a phenomenological Lindblad
bath through the collapse operators
\begin{equation*}
\hat{S}_{\downarrow} = \sqrt{\gamma_{\downarrow}}\,\ket{g}\!\bra{e},
\qquad
\hat{S}_{\uparrow} = \sqrt{\gamma_{\uparrow}}\,\ket{e}\!\bra{g},
\qquad
\hat{S}_{\varphi} = \sqrt{2\gamma_{\varphi}}\,\hat{S}_{e}.
\end{equation*}
Here also a relaxation channel is modelled, where \(\gamma_{\downarrow}\) and \(\gamma_{\uparrow}\) denote downward and
upward population transfer. \(\gamma_{\varphi}\) is the pure-dephasing
rate. In the strict comparison to the monomer benchmark literature, the
zero-temperature limit \(\gamma_{\uparrow}=0\) is used. The corresponding
explicit monomer equations of motion retained for the benchmark implementation
are collected in~\autoref{app:sec:paper_eqs_explicit_eoms}, see in particular
~\cref{eq:app_monomer_eom_sigma_eg,eq:app_monomer_eom_rho_gg}.

\section{Dimer as the \texorpdfstring{$N_{\mathrm{sites}}=2$}{Nsites=2} Realisation}
\label{sec:model_dimer}

The dimer has two optically active sites and is the smallest coupled system
with coherent excitonic mixing and one double-excitation state \(\ket{AB}\).
This is the first system in which cross peaks and waiting-time beats appear.

\subsection{Site Basis and Exciton Basis}

In the site basis \(\{\ket{g},\ket{A},\ket{B},\ket{f}\equiv\ket{AB}\}\), the
field-free Hamiltonian is
\begin{equation*}
\hat{H}_{0}^{(\mathrm{dim})}
=
\omega_{A}\ket{A}\!\bra{A}
+
\omega_{B}\ket{B}\!\bra{B}
+
J\bigl(\ket{A}\!\bra{B}+\ket{B}\!\bra{A}\bigr)
+
(\omega_{A}+\omega_{B})\ket{f}\!\bra{f},
\end{equation*}
where \(\omega_{A}\) and \(\omega_{B}\) are the site energies and \(J\) is the
coherent coupling.

For spectral interpretation, the one-exciton block is diagonalised. The mixing
angle is defined by
\begin{equation}
\tan(2\theta)=\frac{2J}{\omega_{A}-\omega_{B}},
\label{eq:dimer_mixing_angle}
\end{equation}
and the one-exciton eigenenergies are
\begin{equation}
\omega_{2,1}
=
\frac{\omega_{A}+\omega_{B}}{2}
\pm
\frac{1}{2}\sqrt{(\omega_{A}-\omega_{B})^2+4J^2}.
\label{eq:dimer_eigen_energies}
\end{equation}
Generally, the exciton basis is the one used to interpret the results because
the Redfield tensor is evaluated in that basis. The corresponding site-basis and
exciton-basis representation is summarised in~\autoref{fig:rslts_dimer_basis_map}.

\begin{figure}[tbp]
    \centering
    \safeplaceholdergraphic[0.6\linewidth]{figures/rslts_dimer_site_vs_exciton_basis}
    \caption{Basis map used for the dimer model. Left: site basis with local states $\ket{A}$ and $\ket{B}$, site dipoles $\mu_{A}$ and $\mu_{B}$, and coupling $J$. Middle: only the one-exciton block is diagonalised, so $\ket{0}=\ket{g}$ and $\ket{3}=\ket{AB}$ stay unchanged. Right: exciton basis with mixed states $\ket{1}$ and $\ket{2}$ and the adjacent-manifold transitions $\mu_{10}$, $\mu_{20}$, $\mu_{31}$, and $\mu_{32}$.}
    \label{fig:rslts_dimer_basis_map}
\end{figure}

\subsection{Dipole Structure and Excited-State Absorption}

In the exciton basis, the dipole operator becomes
\begin{equation*}
\hat{\mu}^{(\mathrm{dimer})}
=
\sum_{n=A,B}
\Bigl(
\mu_{gn}\,\ket{g}\!\bra{n}
+
\mu_{ng}\,\ket{n}\!\bra{g}
\Bigr)
+
\sum_{n=A,B}
\Bigl(
\mu_{nf}\,\ket{n}\!\bra{f}
+
\mu_{fn}\,\ket{f}\!\bra{n}
\Bigr),
\end{equation*}
with \(\mu_{ng}=\mu_{gn}^{*}\) and \(\mu_{fn}=\mu_{nf}^{*}\). This is also visually represented in~\autoref{fig:rslts_dimer_basis_map}. The
\(\ket{n}\leftrightarrow\ket{f}\) transitions are the additional dipole
channels that distinguish the dimer from the strict monomer benchmark.

Once the system has been driven into the one-exciton manifold, further
absorption into the doubly excited state becomes possible. This is the
excited-state-absorption channel. In the dimer, it is associated with
coherences between \(\ket{f}\) and the one-exciton states and therefore
contributes at detection frequencies different from the ground-to-one-exciton
transitions.

\subsection{Microscopic Bath Coupling}
\label{sec:model_dimer_transfer_model}

For \(N_{\mathrm{sites}}=2\), the site-local coupling operators reduce to
\begin{equation*}
\hat{S}_{A} = \ket{A}\!\bra{A} + \ket{AB}\!\bra{AB},
\qquad
\hat{S}_{B} = \ket{B}\!\bra{B} + \ket{AB}\!\bra{AB}.
\end{equation*}
These operators are inserted into the microscopic interaction Hamiltonian
introduced in~\autoref{eq:oqs_Interaction_Hamiltonian}. In other words, the
dimer bath model used in this thesis is the microscopic site-local bath from
\autoref{chapter:oqs_open_quantum_systems}, specialised to the two dimer sites.
The resulting weak-coupling dynamics is then generated within the Redfield
framework through the tensor expression~\autoref{eq:oqs_Rabcd_full}. For
direct comparison with the benchmark literature, the corresponding explicit
dimer equations of motion are also
presented in~\autoref{app:sec:paper_eqs_explicit_eoms}.
